\documentclass[11pt]{article}
\usepackage[T1]{fontenc}
\usepackage{amsmath,amssymb,amsthm}
\usepackage[margin=1in]{geometry}
\usepackage[hidelinks]{hyperref}
\hypersetup{
  pdftitle={Four- or Eight-Cycles in Twelve-Vertex Two-Attachment Graphs},
  pdfauthor={Xiaojian Zeng},
  pdfsubject={A finite-certificate proof for twelve-vertex two-attachment graphs}
}

\newtheorem{theorem}{Theorem}
\newtheorem{lemma}[theorem]{Lemma}

\title{Four- or Eight-Cycles in Twelve-Vertex\\
Two-Attachment Graphs}
\author{Xiaojian Zeng\\
PBC School of Finance, Tsinghua University, 100190\\
\texttt{xiaojianzeng@foxmail.com}}
\date{}

\begin{document}
\maketitle

\begin{abstract}
We prove that every connected simple graph on twelve vertices with two
distinct vertices of degree two, all other vertices of degree three, and the
two exceptional vertices at distance at least three contains a cycle of
length four or eight.  Adding an edge between the exceptional vertices gives
an exact reduction to marked-edge deletions in the complete census of the 85
connected cubic graphs on twelve vertices.  A certificate covers all 1,110
eligible deletions and supplies a cycle in each.  A separately written
verifier reconstructs the domain, checks every witness, and recomputes its
lexicographic minimality.  The proof depends explicitly on the published
completeness of the cubic census.
\end{abstract}

\section{Statement and context}

A \emph{cycle of length $k$} below is a cyclic sequence of $k$ distinct
vertices with consecutive vertices, including the last and first, adjacent.
The question comes from the study of two-attachment gadgets for graphs without
power-of-two cycle lengths.  Garcia's gadget census treats orders $4,6,8$ and
shows that at order $10$ there is no $C_4/C_8$-free gadget whose attachments
are at distance at least three \cite[Proposition~6.2(2)]{Garcia2026}.
The following theorem settles the corresponding order-$12$ assertion.

\begin{theorem}\label{thm:main}
Let $H$ be a connected simple graph on exactly twelve vertices, and let
$s\ne t$ be vertices such that
\[
 \operatorname{deg}_H(s)=\operatorname{deg}_H(t)=2,\qquad
 \operatorname{deg}_H(v)=3\quad(v\notin\{s,t\}),
 \qquad \operatorname{dist}_H(s,t)\geq 3.
\]
Then $H$ contains a simple cycle of length four or a simple cycle of length
eight.
\end{theorem}

The result is not an immediate consequence of a statement for cubic graphs:
after adding the edge $st$, a cycle in the resulting cubic graph might use
that added edge.  We instead retain the marked edge throughout the finite
reduction.

\section{The marked-edge reduction}

\begin{lemma}\label{lem:bijection}
There is a bijection on labeled objects between graphs $(H,s,t)$ satisfying
the hypotheses of Theorem~\ref{thm:main} and pairs $(G,st)$ such that $G$ is a
connected simple cubic graph on twelve vertices, $st\in E(G)$, $G-st$ is
connected, and
\(\operatorname{dist}_{G-st}(s,t)\geq3\).
The terminal order may be forgotten when testing cycle existence.
\end{lemma}

\begin{proof}
The distance hypothesis implies that $st\notin E(H)$.  Adding $st$ therefore
preserves simplicity and connectedness, and it increases precisely the two
degree-two vertices to degree three.  This produces the required cubic graph
$G$.  Conversely, deleting the marked edge from a pair $(G,st)$ with the
stated connectivity and distance conditions gives exactly the degree sequence
and hypotheses of the theorem.  The two operations are inverse.  Swapping
$s$ and $t$ changes neither the deletion graph nor its cycles, so one test for
each undirected marked edge covers both terminal orientations.
\end{proof}

Meringer's generation method constructs representatives of regular graphs up
to isomorphism by orderly generation and a canonicity test
\cite{Meringer1999}.  The same paper's Table~1 gives 85 connected cubic graphs
on twelve vertices \cite[Table~1]{Meringer1999}, and the associated official
online census provides their adjacency lists \cite{MeringerCensus}.  By
Lemma~\ref{lem:bijection}, it is enough to
test every edge of every representative: an isomorphism from an arbitrary
cubic graph to a representative carries its marked edge to an edge of that
representative, and deletion commutes with the isomorphism.

\section{Finite certificate}

The census has 85 graphs and every cubic graph on twelve vertices has 18
edges, giving 1,530 marked-edge pairs.  The certificate construction found
1,526 connected deletions.  Of these, exactly 1,110 have marked endpoints at
distance at least three.  For each eligible pair it records either a
four-cycle or an eight-cycle that avoids the deleted edge.

For an explicit convention, write a cycle as a vertex tuple beginning at its
least vertex, compare the two orientations lexicographically, and retain the
smaller orientation.  Each certificate record contains the lexicographically
least normalized tuple among all surviving four- and eight-cycles.  This
minimality is not needed for Theorem~\ref{thm:main}, but makes the records
canonical and strengthens reproducibility.

The deterministic verifier performs the following checks from the census and
certificate alone.

\begin{enumerate}
\item It parses all 85 adjacency lists and checks the vertex set, simplicity,
      symmetry of adjacency, degree three, the 18-edge count, and
      connectedness.
\item For every undirected edge, it removes that edge, tests connectivity and
      computes the distance between its endpoints.  It therefore reconstructs
      the eligible set rather than trusting a list supplied by the generator.
\item It rejects duplicate, missing, extra, or ineligible records and checks
      directly that every listed tuple is a simple $C_4$ or $C_8$ in the
      deletion graph.
\item Independently of the generator's vertex-path search, it chooses subsets
      of four and eight surviving edges and recognizes a cycle precisely when
      the chosen subgraph is connected and $2$-regular.  It normalizes all
      cycles found this way and recomputes the least one.
\end{enumerate}

The completed replay returned the counts
\[
\begin{array}{c|r}
\text{census graphs}&85\\
\text{graph--edge pairs}&1{,}530\\
\text{connected deletions}&1{,}526\\
\text{eligible deletions}&1{,}110
\end{array}
\]
and accepted all 1,110 cycles and their claimed minimality.  All predicates
are discrete; the computation uses neither floating-point arithmetic nor
randomization.  Separating discovery from exact replay is also used in other
finite-library results, where independent reconstruction is employed to keep
the generated data outside the verifier's trust boundary
\cite{ZengEtAl2026}; see Section~6 of that paper.

\begin{proof}[Proof of Theorem~\ref{thm:main}]
Given $(H,s,t)$, apply Lemma~\ref{lem:bijection} to obtain a connected simple
cubic graph $G$ with marked edge $st$.  Map $G$ to its representative in the
complete census.  The image of $st$ is one of the 1,110 eligible marked edges,
so the corresponding certificate record gives a $C_4$ or $C_8$ avoiding that
edge.  Pulling the cycle back through the isomorphism gives the same-length
cycle in $G-st=H$.
\end{proof}

\section{Reproducibility and dependency}

The ancillary file \texttt{evidence/cycle\_certificate.json} contains the
records.  From the project directory, the replay command is
\begin{verbatim}
python3 evidence/verify_cycle_certificate.py \
  evidence/external/12_3_3.asc evidence/cycle_certificate.json
\end{verbatim}
The verifier uses only the Python standard library.  It is separately written
from \path{evidence/generate_cycle_certificate.py} and does not import it.

The finite proof is relative to the cited completeness of Meringer's census.
The verifier authenticates and checks the supplied census file and confirms
that its 85 entries have the required graph properties; it does not regenerate
all connected cubic graphs from first principles.  Thus census exhaustion is
an explicit external classification dependency.  No SAT cross-check or proof
assistant is claimed.

\paragraph{Computational and AI assistance.}
Python programs were used to generate and exactly replay the finite
certificate.  OpenAI Codex assisted with literature search, code drafting,
adversarial audit organization, and manuscript preparation.  Language-model
output is not mathematical evidence for the theorem: the argument rests on
Lemma~\ref{lem:bijection}, the cited complete census, the explicit certificate,
and the deterministic verifier.

\bibliographystyle{plain}
\bibliography{references}

\end{document}
